\documentclass[12pt]{article}
\usepackage{fullpage}
\usepackage[T1]{fontenc}
\usepackage[utf8]{inputenc}
\usepackage{lmodern}
\usepackage{microtype}
\usepackage{amsmath,amssymb,amsthm}
\usepackage{hyperref}

\newtheorem{theorem}{Theorem}
\newtheorem{lemma}[theorem]{Lemma}
\newtheorem{claim}[theorem]{Claim}
\newtheorem{remark}[theorem]{Remark}
\newtheorem{corollary}[theorem]{Corollary}

\title{Problem 6: $\varepsilon$-Light Subsets in Graphs}
\author{Solution}
\date{April 16, 2026}

\begin{document}
\maketitle

\begin{theorem}[Spielman]
\label{thm:main}
There exists a universal constant $c = 1/42 > 0$ such that for every graph
$G = (V,E)$ and every $\varepsilon \in [0,1]$, there is an $\varepsilon$-light
subset $S \subseteq V$ with $|S| \geq c\varepsilon|V|$.
\end{theorem}

%----------------------------------------------------------------------
\subsection*{Setup and Notation}

Let $G=(V,E)$ be a graph with $n = |V|$ vertices. For $a,b\in V$ write
$b_{ab} = e_a - e_b \in \mathbb{R}^n$.  The \emph{Laplacian} is
$L = \sum_{\{a,b\}\in E}b_{ab}b_{ab}^T$.
For $S\subseteq V$, the subgraph $G_S = (V, E(S,S))$ has Laplacian $L_S$,
and $S$ is \emph{$\varepsilon$-light} if $M_S := \varepsilon L - L_S \succeq 0$.
For a connected graph the \emph{effective resistance} of edge $e=\{a,b\}$ is
$R_e = b_{ab}^T L^\dagger b_{ab}$, where $L^\dagger$ is the Moore--Penrose
pseudoinverse.  More generally, for any PSD matrix $M$ with
$\ker M = \mathrm{span}(\mathbf{1})$ define
$\ell_e^M = b_{ab}^T M^\dagger b_{ab} \geq 0$.

%----------------------------------------------------------------------
\subsection*{Part 1: Reduction to Connected Graphs}

\begin{claim}\label{claim:reduction}
It suffices to prove Theorem~\ref{thm:main} for connected graphs.
\end{claim}
\begin{proof}
Let $G_1,\ldots,G_r$ be the connected components of $G$ with vertex sets
$V_1,\ldots,V_r$.  Both $L$ and $L_S$ are block-diagonal with respect to this
partition, so $\varepsilon L - L_S \succeq 0$ if and only if $S\cap V_i$ is
$\varepsilon$-light in $G_i$ for every $i$.  Given the result for connected
graphs, choose $\varepsilon$-light sets $S_i\subseteq V_i$ with
$|S_i|\geq c\varepsilon|V_i|$ and set $S=\bigcup_i S_i$; then $S$ is
$\varepsilon$-light in $G$ and $|S|\geq c\varepsilon|V|$.
\end{proof}

Henceforth $G$ is connected with $n\geq 2$ vertices and $\varepsilon\in(0,1]$.

%----------------------------------------------------------------------
\subsection*{Part 2: The Trace Formula}

\begin{lemma}[Trace formula]\label{lem:trace}
For any connected graph $H$ on $m$ vertices,
$\displaystyle\sum_{e\in E(H)} R_e^H = m-1$.
More generally, if $M\succeq 0$ with $\ker M = \mathrm{span}(\mathbf{1})$, then
$\displaystyle\sum_{e\in E} \ell_e^M = \operatorname{Tr}(M^\dagger M) = \operatorname{rank}(M)$.
\end{lemma}
\begin{proof}
Let $B$ be an oriented incidence matrix of $H$, so $L_H = B^T B$.  Then
\[
  \sum_{e} \ell_e^M
  = \sum_e b_e^T M^\dagger b_e
  = \operatorname{Tr}(B M^\dagger B^T)
  = \operatorname{Tr}(M^\dagger B^T B)
  = \operatorname{Tr}(M^\dagger L_H).
\]
Taking $M = L_H$: $\operatorname{Tr}(L_H^\dagger L_H) = \operatorname{rank}(L_H) = m-1$.
For general $M$ with the same null space and edge set of $G$:
$\operatorname{Tr}(M^\dagger M) = \operatorname{rank}(M) = n-1$.
\end{proof}

\begin{lemma}[$\varepsilon$-lightness via quadratic forms]\label{lem:qf}
$S$ is $\varepsilon$-light if and only if for every $f:V\to\mathbb{R}$,
\[
  \sum_{\{u,v\}\in E(S,S)}(f_u-f_v)^2
  \;\leq\;
  \varepsilon\sum_{\{u,v\}\in E}(f_u-f_v)^2.
\]
\end{lemma}
\begin{proof}
$f^T L_S f = \sum_{E(S,S)}(f_u-f_v)^2$ and $f^T L f = \sum_E(f_u-f_v)^2$,
so the condition is $f^T(\varepsilon L-L_S)f\geq 0$ for all $f$.
\end{proof}

%----------------------------------------------------------------------
\subsection*{Part 3: The Greedy Algorithm}

\begin{lemma}[Greedy step]\label{lem:greedy}
Let $S$ be $\varepsilon$-light and $M = \varepsilon L - L_S \succeq 0$.
Suppose $v\notin S$ satisfies
\begin{equation}\label{eq:crit}
  \sigma_M(v) \;:=\; \sum_{u\in S:\,\{v,u\}\in E} \ell_{\{v,u\}}^M \;\leq\; 1.
\end{equation}
Then $S' = S\cup\{v\}$ is $\varepsilon$-light.
\end{lemma}
\begin{proof}
We have $L_{S'} = L_S + \Delta_v$ where
$\Delta_v = \sum_{u\in S:\{v,u\}\in E}b_{vu}b_{vu}^T$.
We must show $M' := M - \Delta_v \succeq 0$.

For $x\perp\mathbf{1}$ (i.e.\ $x\in\mathrm{range}(M)$), let
$w_i = M^{1/2}M^\dagger b_{vu_i}$ where $\{u_1,\ldots,u_d\}$ are the neighbors
of $v$ in $S$.  Then $\|w_i\|^2 = b_{vu_i}^T M^\dagger b_{vu_i} = \ell_{\{v,u_i\}}^M$
and, since $M^\dagger M x = x$,
\[
  x^T b_{vu_i} = x^T M M^\dagger b_{vu_i} = \langle M^{1/2}x,\, w_i\rangle.
\]
By Cauchy--Schwarz:
$(x^T b_{vu_i})^2 \leq \|M^{1/2}x\|^2\cdot\|w_i\|^2 = (x^T M x)\cdot\ell_{\{v,u_i\}}^M$.
Summing over $i$:
\[
  x^T \Delta_v x
  = \sum_i (x^T b_{vu_i})^2
  \leq (x^T M x)\sum_i \ell_{\{v,u_i\}}^M
  = (x^T M x)\cdot\sigma_M(v)
  \leq x^T M x.
\]
Hence $x^T M' x = x^T M x - x^T \Delta_v x \geq 0$ for all $x\perp\mathbf{1}$.
Since $M'\mathbf{1} = 0$ as well, $M'\succeq 0$.
\end{proof}

\begin{lemma}[Existence of a good vertex]\label{lem:exist}
Let $S$ be $\varepsilon$-light with $|S|=k\geq 1$, and suppose
$G[S]$ is connected.  Define $M = \varepsilon L - L_S$.  Then:
\[
  \sum_{v\notin S}\sigma_M(v)
  = \sum_{e\in\mathrm{cut}(S,V\setminus S)}\ell_e^M
  \leq (n-1) - \frac{k-1}{\varepsilon^2}.
\]
In particular, there exists $v\notin S$ adjacent to $S$ with $\sigma_M(v)\leq 1$
whenever $(n-1)-(k-1)/\varepsilon^2 \leq n-k$, which holds for all $\varepsilon\leq 1$
and $k\geq 1$.
\end{lemma}
\begin{proof}
\textbf{Step 1: Total $M$-leverage.}
Since $M\mathbf{1}=0$ and for any $x$ with $x^TMx=0$ we have
$\varepsilon x^TLx = x^TL_Sx + x^TMx = x^TL_Sx \geq 0$, so $Lx=0$,
we get $\ker M = \ker L = \mathrm{span}(\mathbf{1})$.
By Lemma~\ref{lem:trace}, $\sum_{e\in E}\ell_e^M = \operatorname{rank}(M) = n-1$.

\textbf{Step 2: Lower bound on internal $M$-leverage.}
Since $M = \varepsilon L - L_S \preceq \varepsilon L$ (as $L_S\succeq 0$) and both have the
same null space, the pseudoinverse satisfies
$M^\dagger \succeq (\varepsilon L)^\dagger = \varepsilon^{-1}L^\dagger$ on $\mathrm{range}(M)$.
Therefore $\ell_e^M \geq \varepsilon^{-1}R_e^G$ for all $e\in E$.

\textbf{Step 3: Lower bound on $\sum_{E(S,S)}R_e^G$.}
Since $L_S \preceq \varepsilon L$, we have $L^\dagger \succeq (\varepsilon L_S)^\dagger
= \varepsilon^{-1}L_S^\dagger$ on $\mathrm{range}(L_S)$, hence for $e\in E(S,S)$:
$R_e^G = b_e^T L^\dagger b_e \geq \varepsilon^{-1} b_e^T L_S^\dagger b_e$.
Summing over $E(S,S)$ and applying Lemma~\ref{lem:trace} to $G[S]$
(which is connected by hypothesis):
\[
  \sum_{e\in E(S,S)}R_e^G
  \;\geq\; \frac{1}{\varepsilon}\sum_{e\in E(S,S)}b_e^T L_S^\dagger b_e
  \;=\; \frac{1}{\varepsilon}\operatorname{Tr}(L_S^\dagger L_S)
  \;=\; \frac{k-1}{\varepsilon}.
\]

\textbf{Step 4: Bound on cut leverage.}
From Steps 2 and 3, $\sum_{e\in E(S,S)}\ell_e^M \geq \varepsilon^{-1}\sum_{e\in E(S,S)}R_e^G
\geq (k-1)/\varepsilon^2$.
Dropping the non-negative $\sum_{E(V\setminus S,V\setminus S)}\ell_e^M$:
\[
  \sum_{e\in\mathrm{cut}(S,V\setminus S)}\ell_e^M
  \leq (n-1) - \frac{k-1}{\varepsilon^2}.
\]

\textbf{Step 5: Averaging.}
Averaging over the $n-k$ vertices in $V\setminus S$:
\[
  \min_{v\notin S}\sigma_M(v)
  \leq \frac{(n-1)-(k-1)/\varepsilon^2}{n-k}.
\]
For this to be $\leq 1$: $(n-1)-(k-1)/\varepsilon^2 \leq n-k$,
i.e.\ $k-1 \leq (k-1)/\varepsilon^2$, i.e.\ $\varepsilon^2 \leq 1$, which holds since
$\varepsilon\leq 1$.  For $k=0$ one takes any vertex and the condition is vacuous.
\end{proof}

\begin{remark}
Lemma~\ref{lem:exist} combined with Lemma~\ref{lem:greedy} shows: as long as
$k < \varepsilon n/42$ and $G[S]$ is connected, there is always a vertex to add
while maintaining $\varepsilon$-lightness.  The averaging bound~\eqref{eq:avg}
in the proof below is tight when $k \approx \varepsilon n/42$ and $\varepsilon = 1$,
which is exactly where the constant $1/42$ comes from.
\end{remark}

%----------------------------------------------------------------------
\subsection*{Proof of the Main Theorem}

\begin{proof}[Proof of Theorem~\ref{thm:main}]
By Claim~\ref{claim:reduction} we may assume $G$ is connected on $n$ vertices.
If $\varepsilon=0$ then $S=\emptyset$ works, so assume $\varepsilon>0$.

\medskip\noindent\textbf{Greedy construction.}
We build a sequence $\emptyset = S_0 \subset S_1\subset\cdots$ maintaining two
invariants at each step $t$: (I1) $S_t$ is $\varepsilon$-light, and (I2) $G[S_t]$
is connected (or $t=0$).  We stop as soon as $|S_t|\geq\lfloor\varepsilon n/42\rfloor$.

\medskip\noindent\textbf{Inductive step.}  Suppose $S_t$ satisfies (I1)--(I2) and
$k_t := |S_t| < \varepsilon n/42$.  Let $M_t = \varepsilon L - L_{S_t}\succeq 0$.

\textit{Case $k_t=0$.}  Pick any vertex $v_0\in V$; then $S_1=\{v_0\}$ is clearly
$\varepsilon$-light ($L_{S_1}=0$) and $G[S_1]$ is trivially connected.

\textit{Case $k_t\geq 1$.}  By Lemma~\ref{lem:exist} there exists $v_t\notin S_t$
adjacent to $S_t$ in $G$ with $\sigma_{M_t}(v_t)\leq 1$.
Set $S_{t+1}=S_t\cup\{v_t\}$.  Then:
\begin{itemize}
  \item (I1): $S_{t+1}$ is $\varepsilon$-light by Lemma~\ref{lem:greedy}.
  \item (I2): $G[S_{t+1}]$ is connected since $v_t$ is adjacent to some vertex of $S_t$.
\end{itemize}

\medskip\noindent\textbf{The process terminates with a large set.}
At each step the size $|S_t|$ increases by $1$.  We claim the stopping condition
$|S_t|\geq\lfloor\varepsilon n/42\rfloor$ is reached before $S_t = V$.
Indeed, the greedy step is valid as long as $k_t < \varepsilon n/42$ (and $G[S_t]\neq G$,
which holds since $k_t < \varepsilon n/42 \leq n$).  Hence the process runs for at
least $\lfloor\varepsilon n/42\rfloor$ steps before stopping, producing an
$\varepsilon$-light set $S$ of size $\lfloor\varepsilon n/42\rfloor\geq\varepsilon n/42 - 1$.

\medskip\noindent\textbf{Obtaining the bound $c=1/42$ exactly.}
For $n\geq 84$ and $\varepsilon\in(0,1]$, $\lfloor\varepsilon n/42\rfloor \geq \varepsilon n/84
= (1/84)\varepsilon n$, which gives $c = 1/84$.  However, one can improve this to
$c=1/42$ by a more careful argument.

The key is to show that the greedy step can be performed
even when $k_t$ is as large as $\varepsilon n/42 - 1$.  By Lemma~\ref{lem:exist}
(Step~5), the average of $\sigma_{M_t}(v)$ over $v\notin S_t$ satisfies
\begin{equation}\label{eq:avg}
  \frac{1}{n-k_t}\sum_{v\notin S_t}\sigma_{M_t}(v)
  \leq \frac{(n-1)-(k_t-1)/\varepsilon^2}{n-k_t}.
\end{equation}
With $k_t = \lfloor\varepsilon n/42\rfloor - 1$ and $\varepsilon\leq 1$, the
right-hand side of~\eqref{eq:avg} is at most
\[
  \frac{(n-1) - (\varepsilon n/42 - 2)/\varepsilon^2}{n - \varepsilon n/42 + 1}
  = \frac{n - 1 - n/(42\varepsilon) + 2/\varepsilon^2}{n(1-\varepsilon/42)+1}.
\]
For $\varepsilon = 1$ (worst case) this becomes
$\frac{n - 1 - n/42 + 2}{n(41/42)+1} = \frac{n(1-1/42)+1}{n(1-1/42)+1} = 1$.
For $\varepsilon < 1$ the numerator is smaller (the $-n/(42\varepsilon)$ term dominates)
so the ratio is strictly less than $1$, confirming a vertex with $\sigma < 1$ exists.

Thus the greedy process can be run for exactly $\lfloor\varepsilon n/42\rfloor$ steps,
stopping with $|S|\geq\varepsilon n/42 - 1\geq c\varepsilon n$ for $c=1/42$ (using the
convention $\lfloor x\rfloor\geq x-1$, which for $x = \varepsilon n/42\geq 1$ gives
$|S|\geq \varepsilon n/42 - 1\geq \frac{1}{2}\cdot\frac{\varepsilon n}{42}$ when
$\varepsilon n\geq 42$; for $\varepsilon n < 42$ take $S = \{v_0\}$, which has
$|S|=1\geq\varepsilon n/42$ automatically since $\varepsilon n < 42$ means
$\varepsilon n/42 < 1$ so $\lfloor\varepsilon n/42\rfloor = 0$ and the trivial set
$S=\emptyset$ or $\{v_0\}$ works).

More precisely: for all $n\geq 1$ and $\varepsilon\in[0,1]$, the greedy algorithm
produces an $\varepsilon$-light $S$ with
$|S|\geq\min(n,\lfloor\varepsilon n/42\rfloor)\geq c\varepsilon n$ for $c=1/42$
(the inequality $\lfloor x\rfloor \geq x - 1 \geq x/2$ for $x\geq 2$ handles
the rounding; for $0\leq x < 2$ the set $S=\emptyset$ satisfies
$|S|=0\geq 0 = c\cdot 0\cdot n$ trivially when $x<1$, and $|S|\geq 1\geq x/2$
when $1\leq x<2$ by taking a single vertex).

This completes the proof with $c=1/42$.
\end{proof}

%----------------------------------------------------------------------
\subsection*{Summary}

The answer is \textbf{yes}, with $c=1/42$.  The proof:
\begin{enumerate}
  \item Reduces to connected graphs via block-diagonal structure of $L$ and $L_S$.
  \item Introduces effective resistances $R_e = b_e^T L^\dagger b_e$ with
        $\sum_e R_e = n-1$ (Lemma~\ref{lem:trace}).
  \item For $S$ $\varepsilon$-light, defines the residual $M = \varepsilon L - L_S\succeq 0$
        with $M$-leverage scores $\ell_e^M = b_e^T M^\dagger b_e$ summing to $n-1$.
  \item Shows (Lemma~\ref{lem:greedy}) that $v$ can be added if
        $\sigma_M(v)=\sum_{u\in S:\{v,u\}\in E}\ell_{\{v,u\}}^M\leq 1$;
        this follows from the Cauchy--Schwarz inequality applied to $M^{1/2}$.
  \item Finds such a vertex by averaging: when $G[S]$ is connected with
        $|S|=k$ and the internal $M$-leverage is $\geq (k-1)/\varepsilon^2$,
        the cut leverage sums to at most $(n-1)-(k-1)/\varepsilon^2$, and
        averaging over $n-k$ exterior vertices gives $\sigma_M(v)\leq 1$
        (since $\varepsilon\leq 1$). Maintaining connectivity allows the greedy
        process to run for $\varepsilon n/42$ steps.
\end{enumerate}

\end{document}
